\chapter{Designing a protected random observation clock}
\label{rr:clock}

Timing is part of experimental access. A fixed-time population comparison can
have either sign in a classical network, while its uniform average starting at
preparation has a protected sign. This chapter describes clocks that retain
that guarantee and makes a limited resource-design problem exactly solvable.
The mixture representation and linear-programming geometry are established
mathematical tools; their use here does not assert new generic optimization
theory.

\section{Decreasing densities are mixtures of uniform clocks}
For a protected classical reset state $j$ and any other initial distribution
$\rho$, write
\[
 f(t)=P_j(X_t=j)-P_\rho(X_t=j),\qquad
 H(T)=\int_0^T f(t)\,dt\geq0.
\]
The same statement applies to a region with the required regenerative reset
property. A clock is drawn independently of the physical trajectory, with a
common probability density $w$ for both preparations. Its contrast is
$C_w=\int_0^\infty w(t)f(t)\,dt$.

\begin{proposition}[A constructive protected clock class]
Every nonincreasing probability density on $[0,\infty)$ is, up to changes on
a set of measure zero, a mixture of uniform densities starting at zero:
\[
 w(t)=\int_{(t,\infty)}\frac{1}{T}\,\mu(dT),\qquad
 \mu\ \hbox{a probability measure on }(0,\infty).
\]
Consequently $C_w=\int H(T)/T\,\mu(dT)\geq0$. When the expected duration
exists, $\mathbb E_w t=\frac12\int T\,\mu(dT)$.
\end{proposition}
\begin{proof}
Take a right-continuous nonincreasing representative and its nonnegative
Stieltjes measure $-dw$. Set $\mu(dT)=-T\,dw(T)$. Integrability and
monotonicity imply $T w(T)\to0$ at both endpoints. Integration by parts gives
$\int\mu(dT)=\int_0^\infty w(t)\,dt=1$, and Stieltjes integration gives
the displayed representation. Since $|f|\leq1$, Fubini applies to the contrast;
each mixture component contributes $H(T)/T$. The duration formula follows
from the mean of a uniform distribution.
\end{proof}

The construction gives a direct implementation: choose a horizon $T$ using
$\mu$, then draw the actual readout time uniformly from $[0,T]$. The horizon
mixture and readout time are chosen independently of the current system
trajectory. A common exponential clock belongs to this class, but some clocks
erase a particular quantum signal. Classical validity and signal strength
are separate design questions.

This is a sufficient clock class, not a characterization of every density
valid for every classical model. A shifted uniform window or an increasing
density need not preserve the guarantee. The independent tests supply fully
classical counterexamples for both.

\section{Exact design on a finite menu}
For each available uniform horizon $T_k$, suppose a model calculation supplies
an enclosing signed-contrast interval $[l_k,u_k]$. The desired reversal
advantage is $v_k=-C_k$. Let $c_k\geq0$ be a calibrated expected acquisition
cost and $B$ an expected-cost budget. A randomized horizon choice has weights
$p_k\geq0$, $\sum_kp_k=1$. Its guaranteed advantage is at least
$\sum_kp_k(-u_k)$. The finite design problem is
\begin{equation}\label{rr:clock-lp}
 \max_p\sum_kp_k v_k\quad\hbox{subject to}\quad
 \sum_kp_kc_k\leq B,\quad \sum_kp_k=1,\quad p_k\geq0,
\end{equation}
first with $v_k=-u_k$ and then with $v_k=-l_k$.

\begin{proposition}[Two horizons and a checkable global certificate]
If the menu is finite and feasible, problem~\eqref{rr:clock-lp} has an optimum
supported on at most two horizons. An admissible $p$, a slope $a\geq0$ and an
intercept $b$ certify global optimality if
\[
 b+a c_k\geq v_k\quad\hbox{for every }k,
 \qquad\sum_kp_kv_k=b+aB.
\]
\end{proposition}
\begin{proof}
An extreme feasible point has at most two positive coordinates: with three,
a nonzero supported perturbation can preserve both total probability and
total cost, giving a feasible line segment through the point. If the cost
constraint is slack, an extreme point has one positive coordinate. A linear
objective attains an optimum at an extreme point of this compact polytope.
For any feasible $p'$, the supporting inequalities give
$\sum_kp'_kv_k\leq b+a\sum_kp'_kc_k\leq b+aB$. Equality at $p$ proves
optimality without repeating the optimization.
\end{proof}

The implementation enumerates feasible singletons and pairs, then exports
the supporting line. A separate checker verifies feasibility and the dual
inequalities using exact rational arithmetic. Solving with $-u_k$ gives an
achievable guaranteed advantage; solving with $-l_k$ gives an upper bound on
the unknown true optimum. Their difference is an explicit numerical-enclosure
gap. An interval covering negative advantage is not declared a successful
reversal experiment.

The quantum example uses sixteen horizons $T=k/4$, $k=1,\ldots,16$,
depolarization rate $1/10$, rotation frequency $2$, and detector gain $9/10$.
Exact enclosures are rounded outward, never to the nearest favorable value.
The supplied budgets constrain mean active evolution time $T/2$. A physical
energy model must additionally account for preparation, control, measurement
and any duration-dependent supplied power. The code accepts such calibrated
costs but does not infer them from the horizon labels.

\section{When mixing does not improve an energy objective}
A timing mixture can improve a constrained tradeoff while providing no gain
for a different objective. Suppose $c_k>0$ and the chosen sufficient precision
bound depends on advantage squared per expected energy. For any mixture with
positive average advantage,
\[
 \frac{(\sum_kp_kv_k)_+^2}{\sum_kp_kc_k}
 \leq\sum_kp_k\frac{(v_k)_+^2}{c_k}
 \leq\max_k\frac{(v_k)_+^2}{c_k}.
\]
Here $x_+=\max\{0,x\}$. Replace $v_k$ by its positive part and apply
Cauchy--Schwarz to $\sqrt{p_kc_k}$ and
$\sqrt{p_k}(v_k)_+/\sqrt{c_k}$. Thus a mixture cannot beat the best single
horizon for this particular proxy. Exact binomial power, preparation-dependent
overheads, simultaneous objectives or additional constraints require their
own calculation. Neither this proxy nor the mean-cost budget supplies a
pathwise expenditure guarantee for a stopped sequential experiment.

\section{Clock calibration and adaptive use}
If the actual common clock law differs from its declared protected law by
total variation at most $\zeta$, then its contrast can change by at most
$2\zeta$, because $f\in[-1,1]$. If each arm has its own discrepancy radius
$\zeta_j,\zeta_\rho$, a bound is $\zeta_j+\zeta_\rho$, since each
population lies in $[0,1]$. A fixed nonnegative readout delay $d$ shifts a
uniform $[0,T]$ clock; the total variation from the original law is
$\min\{1,d/T\}$. Delay therefore belongs in the physical model or the
calibration budget. It cannot be dismissed as an innocuous change of origin.

For sequential testing, select the horizon mixture and betting rule using
completed earlier trials. The current arm and readout time must then be
drawn without allowing those fresh values to select the bet. The classical
null is valid after averaging the clock; it need not be valid conditional on
its realized time. Later trials may use all earlier timing records. This
distinction permits adaptive design while preventing selection of the very
time intervals in which a classical instantaneous contrast is negative.

\begin{verbatim}
from constructive_access.robust import clock_design

certificate = clock_design.design(menu, budget)
assert clock_design.verify(menu, budget, certificate)
\end{verbatim}
The menu is a supplied physical-model and calibration contract. The checker
certifies the declared design problem, not the truth of the apparatus model.
